\section{ماژول uniql\_compressor (UniQL Compressor)}
\label{sec:module_uniql_compressor}

\subsection{معرفی و انگیزه}
این ماژول در چارچوب پژوهش حاضر پیاده‌سازی و ارزیابی شده و فاقد مقاله مستقل جداگانه در ادبیات تحقیق است.

ترکیب فشرده‌سازی کم‌مرتبه (\lr{Low-Rank Decomposition}) و کوانتیزاسیون گسسته (\lr{Quantization}) یکی از کارآمدترین روش‌ها برای فشرده‌سازی مضاعف شبکه‌های عصبی عمیق به شمار می‌رود. در حالی که روش‌های تجزیه کم‌رتبه ماتریس‌ها را به فاکتورهای کوچکتری تجزیه می‌کنند، مقادیر عددی این فاکتورها همچنان در فرمت ممیز شناور ۳۲ بیتی یا ۱۶ بیتی نگهداری می‌شوند.

موتور فشرده‌سازی یکپارچه \lr{UniQL Compressor (Joint Low-Rank SVD and Uniform Quantization)} به منظور تجمیع همزمان تقریب کم‌مرتبه و کوانتیزاسیون یکنواخت بازه‌ای در یک مرحله واحد توسعه یافته است. در این الگوریتم، ابتدا نگاشت‌های خطی بزرگ با بهره‌گیری از تجزیه مقادیر منفرد (\lr{Truncated SVD}) به دو زیرفاکتور کم‌رتبه افراز شده و بلافاصله بر روی ماتریس بازسازی‌شده یا فاکتورهای تقریب، نگاشت کوانتیزاسیون یکنواخت با تنظیم تطبیقی بیت اعمال می‌شود تا از واگرایی خطای بازسازی جلوگیری به عمل آید.

\subsection{مبانی نظری و فرمول‌بندی ریاضی}
فرض کنید ماتریس وزن اولیه یک لایه به صورت $\mathbf{W} \in \mathbb{R}^{d_{\text{out}} \times d_{\text{in}}}$ باشد. در گام نخست، تجزیه مقادیر منفرد بر روی ماتریس وزن صورت می‌گیرد:
\begin{equation}
\mathbf{W} = \mathbf{U} \mathbf{\Sigma} \mathbf{V}^\top
\end{equation}
با انتخاب رتبه هدف $r \ll \min(d_{\text{out}}, d_{\text{in}})$، تقریب کم‌مرتبه مرتبه $r$ به صورت فاکتورهای متقارن بازسازی می‌شود:
\begin{equation}
\mathbf{A} = \mathbf{U}_{:, :r} \mathbf{\Sigma}_{:r}^{1/2} \in \mathbb{R}^{d_{\text{out}} \times r}, \quad \mathbf{B} = \mathbf{\Sigma}_{:r}^{1/2} \mathbf{V}_{:, :r}^\top \in \mathbb{R}^{r \times d_{\text{in}}}
\end{equation}
ماتریس حاصل‌ضرب بازسازی‌شده کم‌رتبه برابر با $\hat{\mathbf{W}}_{\text{lr}} = \mathbf{A} \mathbf{B}$ خواهد بود. در گام دوم، نگاشت کوانتیزاسیون یکنواخت با تعداد بیت مشخص $b$ (به طور پیش‌فرض $b=8$) بر روی تانسور وزن اعمال می‌گردد. عملگر کوانتیزاسیون یکنواخت کمینه-بیشینه (\lr{Uniform Min-Max Quantization}) به صورت زیر تعریف می‌شود:
\begin{equation}
q_{\text{max}} = 2^b - 1, \quad s = \frac{\max(\mathbf{W}) - \min(\mathbf{W})}{q_{\text{max}}}
\end{equation}
\begin{equation}
\mathcal{Q}(\mathbf{W}, b) = \operatorname{clip}\left(\left\lfloor \frac{\mathbf{W} - \min(\mathbf{W})}{s} \right\rceil, 0, q_{\text{max}}\right) \cdot s + \min(\mathbf{W})
\end{equation}
سپس خطای نسبی بازسازی فروبنیوس مورد پایش قرار می‌گیرد:
\begin{equation}
\mathcal{E}_{\text{rel}} = \frac{\|\mathbf{W} - \mathcal{Q}(\hat{\mathbf{W}}_{\text{lr}}, b)\|_F}{\|\mathbf{W}\|_F + \epsilon}
\end{equation}
در صورتی که خطای نسبی بازسازی از آستانه مجاز $\mathcal{E}_{\text{max}}$ (به عنوان مثال $0.20$) فراتر رود، فرآیند تلاش مجدد تطبیقی با افزایش عمق بیتی ($b \leftarrow b+1$) و افزایش رتبه ($r \leftarrow r \times 1.25$) آغاز شده و در صورت عدم همگرایی، لایه به حالت امن اولیه بازگردانده می‌شود (\lr{Fallback}).

\subsection{معماری و پیاده‌سازی}
کلاس موتور فشرده‌سازی \lr{UniQLCompressor} در فایل \lr{neuralops/compression\_engine/uniql\_compressor.py} مستقر بوده و تحت نام‌های \lr{uniql} و \lr{uniql\_compressor} در رجیستری مرکزی بارگذاری می‌شود. شاخص‌ترین قابلیت‌های طراحی این موتور عبارتند از:
\begin{itemize}
    \item \textbf{حلقه تنظیم مجدد تطبیقی (\lr{Adaptive Bit \& Rank Retries}):} اجرای فرآیند تکرار شونده تا سقف \lr{max\_retries=3} که در صورت بالا بودن خطای تقریب اولیه، به صورت هوشمند و خودکار نرخ بیت یا رتبه را افزایش می‌دهد تا بازنمایی معنایی لایه حفظ گردد.
    \item \textbf{پشتیبانی ساختاری از لایه‌های ترانهاده \lr{Conv1D}:} شناسایی خودکار ویژگی‌های لایه‌های کانولوشن یک‌بعدی در مدل‌های دیکودر (\lr{GPT-2}) و اعمال مناسب ترانهاده قبل و بعد از تجزیه.
    \item \textbf{محافظت از ساختارهای بحرانی شبکه:} مستثنی کردن خودکار لایه‌های پایانی سر طبقه‌بندی، واژگان و تعبیه از فشرده‌سازی به منظور تثبیت احتمالات خروجی.
    \item \textbf{مسیر امن بازگشت به وزن‌های اصلی (\lr{Safe Fallback}):} جایگزینی لایه اصلی دست‌نخورده در صورتی که خطای نهایی پس از تمامی تلاش‌ها فراتر از آستانه اطمینان باشد.
\end{itemize}

\subsection{نتایج تجربی}
ارزیابی بنچ‌مارک فشرده‌ساز \lr{uniql\_compressor} بر روی مدل‌های زبانی در جدول \ref{tab:benchmark_uniql_compressor} ارائه گردیده است:

\begin{table}[H]
\centering
\small
\caption{نتایج تجربی و اعتبارسنجی بنچ‌مارک فشرده‌ساز \lr{uniql\_compressor} روی مدل‌های پایه}
\label{tab:benchmark_uniql_compressor}
\begin{tabular}{lccccc}
\toprule
\textbf{مدل پایه} \lr{(Base Model)} & \textbf{دقت} \lr{(Accuracy)} & \textbf{شاخص} \lr{(F1-Score)} & \textbf{حافظه} \lr{(Peak VRAM)} & \textbf{توان} \lr{(Throughput)} & \textbf{وضعیت} \\
\midrule
\lr{DistilBERT} & \lr{90.14\%} & \lr{0.9027} & \lr{339.22 MB} & \lr{150,882.16 tok/s} & \lr{PASS} \\
\lr{RoBERTa} & \lr{82.45\%} & \lr{0.8468} & \lr{359.04 MB} & \lr{104,803.31 tok/s} & \lr{PASS} \\
\lr{GPT-2} & \lr{60.89\%} & \lr{0.6634} & \lr{635.04 MB} & \lr{72,184.34 tok/s} & \lr{PASS} \\
\lr{Qwen} & \lr{---} & \lr{---} & \lr{---} & \lr{---} & \lr{NOT\_RUN} \\
\bottomrule
\end{tabular}
\end{table}

بر پایه داده‌های اعتبارسنجی بنچ‌مارک، ماژول \lr{uniql\_compressor} به یکی از درخشان‌ترین نتایج فشرده‌سازی در کل سیستم دست یافته است. در مدل \lr{DistilBERT} با ثبت دقت قابل توجه \lr{90.14\%} و نرخ پردازش بسیار بالای \lr{150,882.16} توکن بر ثانیه، پایداری کامل مدل تضمین شده است. همچنین در مدل \lr{RoBERTa} دقت \lr{82.45\%} با شاخص F1 برابر با \lr{0.8468} و در مدل تولیدی \lr{GPT-2} دقت \lr{60.89\%} با وضعیت \lr{PASS} حاصل گردیده است. مدل \lr{Qwen} در ارزیابی جامع یکپارچه اجرا نشده و به صورت \lr{NOT\_RUN} نشانه‌گذاری شده است.

\subsection{محدودیت‌ها و تحلیل شکست}
با وجود نتایج تجربی برجسته، محدودیت‌های تحلیلی این روش به شرح زیر است:
\begin{itemize}
    \item \textbf{شبیه‌سازی کوانتیزاسیون به جای فرمت سخت‌افزاری واقعی (\lr{Fake Quantization}):} پیاده‌سازی فعلی عملگر کوانتیزاسیون را با گردکردن مقادیر در تانسورهای ممیز شناور پیاده می‌سازد؛ بنابراین کاهش واقعی فیزیکی مصرف حافظه در سطح ذخیره‌سازی نوع داده به فرمت عدد صحیح ۸ بیتی سخت‌افزاری (\lr{INT8 INT4 Tensor Cores}) نیازمند الحاق هسته‌های محاسباتی سفارشی است.
    \item \textbf{خطای تجمیعی در شبکه‌های دارای مقادیر پرت شدید:} در صورتی که دامنه‌ی مقادیر وزن در لایه‌ای بسیار گسترده باشد، کوانتیزاسیون یکنواخت مقادیر کوچک را با تقریب درشت گرد می‌کند که می‌تواند منجر به فعال‌سازی مکرر حلقه تلاش مجدد و انتخاب رتبه‌های بالاتر گردد.
\end{itemize}
